\chapter{Electrical and Electronic Design}\label{ch:electrical-and-electronic-design}

\section{Electrical Architecture}\label{electrical-architecture}

The electrical system has two energy domains. The high-current domain
supplies the four actuators from the six-cell lithium-polymer battery.
The logic domain supplies the controller, IMU, transceiver, receiver,
and display through regulated rails. Separating these functions in the
schematic and layout makes it possible to review conductor ratings,
protection, startup behavior, and noise coupling independently while
maintaining one intentional system reference.

\begin{longtable}[]{@{}
  >{\raggedright\arraybackslash}p{(\columnwidth - 0\tabcolsep) * \real{1.0000}}@{}}
\toprule\noalign{}
\begin{minipage}[b]{\linewidth}\raggedright
\textless INSERT FIGURE 5.1: Show VBAT, protected or switched battery
output, four actuator branches, logic converter input, 5 V rail, 3.3 V
rail, controller, IMU, CAN transceiver, receiver, display, debug
connection, and every ground-return path. Use distinct line weights for
high-current and logic paths\textgreater{}\\
What the figure should show: Show VBAT, protected or switched battery
output, four actuator branches, logic converter input, 5 V rail, 3.3 V
rail, controller, IMU, CAN transceiver, receiver, display, debug
connection, and every ground-return path. Use distinct line weights for
high-current and logic paths\\
Likely source: final KiCad hierarchical schematic\strut
\end{minipage} \\
\midrule\noalign{}
\endhead
\bottomrule\noalign{}
\endlastfoot
\end{longtable}

\begin{center}\singlespacing
\phantomsection\addcontentsline{lof}{figure}{\protect\numberline{5.1}Electrical architecture and voltage domains.}
\textbf{Figure 5.1. Electrical architecture and voltage domains.}
\end{center}

\par\medskip\noindent\singlespacing
\phantomsection\addcontentsline{lot}{table}{\protect\numberline{5.1}Electrical components and implementation status}
\textbf{Table 5.1. Electrical components and implementation status}\par

\begin{longtable}[]{@{}
  >{\raggedright\arraybackslash}p{(\columnwidth - 6\tabcolsep) * \real{0.1680}}
  >{\raggedright\arraybackslash}p{(\columnwidth - 6\tabcolsep) * \real{0.2480}}
  >{\raggedright\arraybackslash}p{(\columnwidth - 6\tabcolsep) * \real{0.2480}}
  >{\raggedright\arraybackslash}p{(\columnwidth - 6\tabcolsep) * \real{0.3360}}@{}}
\toprule\noalign{}
\begin{minipage}[b]{\linewidth}\raggedright
\textbf{Function}
\end{minipage} & \begin{minipage}[b]{\linewidth}\raggedright
\textbf{Selected or proposed component}
\end{minipage} & \begin{minipage}[b]{\linewidth}\raggedright
\textbf{Evidence\slash{}status}
\end{minipage} & \begin{minipage}[b]{\linewidth}\raggedright
\textbf{Design issue to close}
\end{minipage} \\
\midrule\noalign{}
\endhead
\bottomrule\noalign{}
\endlastfoot
Main controller & ESP32 NodeMCU prototype & Active firmware platform &
Exact board revision, regulator capacity, final PCB integration \\
Future controller & NUCLEO-H723ZG / STM32H723ZG & Scaffolded migration
target & Migration decision, pin freeze, timing verification \\
IMU & ST ISM330DHCX & SPI driver and estimator path present & Final
mounting, axes, supply, interrupts, calibration \\
CAN physical layer & Adafruit CAN Pal using TJA1051T\slash{}3 & Prototype
interface & Final footprint or direct transceiver, standby pin,
protection, termination \\
Actuators & 4 \ensuremath{\times} CubeMars AK40-10 & Motor communication\slash{}motion partially
demonstrated & IDs, modes, limits, firmware, harness pinout \\
Radio receiver & RadioMaster RP2 V2 ExpressLRS & CRSF decoding reported
& Supply, UART mapping, link-loss behavior \\
Display & Waveshare 2-inch ST7789 module & UI interface reported & Final
3.3 V compatibility, connector footprint, update load \\
5 V conversion & LT8609S design & Schematic reviewed; hardware
validation unavailable & Enable threshold, layout, thermal and transient
testing \\
Battery protection & BQ76942-based 6S design discussed &
Recommended\slash{}separate BMS work; not confirmed manufactured & Final
architecture, MOSFETs, shunt, cell connections, configuration \\
\end{longtable}

\section{Battery, Protection, and Power Distribution}\label{battery-protection-and-power-distribution}

A six-cell lithium-polymer pack has a nominal voltage of approximately
22.2 V and a fully charged voltage of 25.2 V. Converter and protection
ratings must use the maximum voltage plus credible transients, not the
nominal label. The selected actuators are rated for 24 V operation
\cite{cubemars2026ak4010}, so the design supplies them from the battery domain rather
than generating a regulated 24 V rail. The exact battery capacity,
discharge rating, connector, cell condition, and allowable undervoltage
threshold remain hardware verification items.

The high-current design should include a master disconnect, an emergency
method that removes actuator energy, branch overcurrent protection,
reverse-polarity prevention or keyed connectors, and a precharge or
anti-spark strategy if connector inrush is material. A BQ76942 can
monitor and protect three to ten series cells \cite{ti2026bq76942}, but using the
device requires a complete cell-tap network, current shunt, charge and
discharge protection MOSFETs, thermistors, pack\slash{}load connections,
configuration, and fault testing. The thesis must distinguish a separate
BMS board from the main controller board if they are manufactured
independently.

\begin{longtable}[]{@{}
  >{\raggedright\arraybackslash}p{(\columnwidth - 0\tabcolsep) * \real{1.0000}}@{}}
\toprule\noalign{}
\begin{minipage}[b]{\linewidth}\raggedright
\textless INSERT FIGURE 5.2: Show the exact 6S battery connector or two
connector branches, master switch or contactor, emergency stop,
precharge or anti-spark element, pack fuse, per-motor or per-pair fuses,
BMS\slash{}protection boundaries, current shunt, actuator connectors, LT8609S
input, and high-current returns. Label ratings and connector pin
numbers\textgreater{}\\
What the figure should show: Show the exact 6S battery connector or two
connector branches, master switch or contactor, emergency stop,
precharge or anti-spark element, pack fuse, per-motor or per-pair fuses,
BMS\slash{}protection boundaries, current shunt, actuator connectors, LT8609S
input, and high-current returns. Label ratings and connector pin
numbers\\
Likely source: final power schematic and wiring harness drawing\strut
\end{minipage} \\
\midrule\noalign{}
\endhead
\bottomrule\noalign{}
\endlastfoot
\end{longtable}

\begin{center}\singlespacing
\phantomsection\addcontentsline{lof}{figure}{\protect\numberline{5.2}Battery and high-current power-distribution schematic.}
\textbf{Figure 5.2. Battery and high-current power-distribution schematic.}
\end{center}

\clearpage
\begin{landscape}
\par\medskip\noindent\singlespacing
\phantomsection\addcontentsline{lot}{table}{\protect\numberline{5.2}Electrical power budget}
\textbf{Table 5.2. Electrical power budget}\par

\begin{longtable}[]{@{}
  >{\raggedright\arraybackslash}p{(\columnwidth - 10\tabcolsep) * \real{0.1812}}
  >{\raggedright\arraybackslash}p{(\columnwidth - 10\tabcolsep) * \real{0.0872}}
  >{\raggedright\arraybackslash}p{(\columnwidth - 10\tabcolsep) * \real{0.1342}}
  >{\raggedright\arraybackslash}p{(\columnwidth - 10\tabcolsep) * \real{0.2081}}
  >{\raggedright\arraybackslash}p{(\columnwidth - 10\tabcolsep) * \real{0.2081}}
  >{\raggedright\arraybackslash}p{(\columnwidth - 10\tabcolsep) * \real{0.1812}}@{}}
\toprule\noalign{}
\begin{minipage}[b]{\linewidth}\raggedright
\textbf{Load}
\end{minipage} & \begin{minipage}[b]{\linewidth}\raggedright
\textbf{Quantity}
\end{minipage} & \begin{minipage}[b]{\linewidth}\raggedright
\textbf{Voltage}
\end{minipage} & \begin{minipage}[b]{\linewidth}\raggedright
\textbf{Rated\slash{}expected current}
\end{minipage} & \begin{minipage}[b]{\linewidth}\raggedright
\textbf{Peak\slash{}design current}
\end{minipage} & \begin{minipage}[b]{\linewidth}\raggedright
\textbf{Evidence}
\end{minipage} \\
\midrule\noalign{}
\endhead
\bottomrule\noalign{}
\endlastfoot
AK40-10 actuator & 4 & VBAT & 2.7 A each manufacturer rated current &
7.3 A each manufacturer peak current & Manufacturer specification;
chosen limits unverified \\
ESP32 controller board & 1 & 5 V input or 3.3 V rail &
\textless MEASURE\textgreater{} & \textless MEASURE startup\slash{}transmit
peak\textgreater{} & No final measured data \\
ISM330DHCX & 1 & 3.3 V & \textless CALCULATE FROM SELECTED
MODE\textgreater{} & \textless VERIFY\textgreater{} & Use exact
datasheet operating mode \cite{st2026ism330dhcx} \\
CAN transceiver & 1 & 3.3 V module input &
\textless VERIFY\textgreater{} & \textless VERIFY dominant-bus
load\textgreater{} & Adafruit\slash{}NXP documentation \cite{nxp2026tja1051}, \cite{adafruit2026canpal} \\
Radio receiver & 1 & 5 V & \textless MEASURE\textgreater{} &
\textless MEASURE\textgreater{} & Confirm installed RP2 V2 revision
\cite{radiomaster2026manuals} \\
ST7789 display module & 1 & 3.3 V & \textless MEASURE at selected
brightness\textgreater{} & \textless MEASURE\textgreater{} & Confirm
module revision \cite{waveshare2026lcd} \\
Logic total & --- & 5 V and 3.3 V & \textless SUM WITH
MARGIN\textgreater{} & \textless LOAD-STEP TEST\textgreater{} &
Pending \\
\end{longtable}
\end{landscape}
\clearpage

\emph{Note.} Manufacturer peak ratings are not simultaneous system
requirements. Final branch and pack ratings must follow the configured
motor limits and measured duty cycle.

\section{Voltage Conversion and Logic Rails}\label{voltage-conversion-and-logic-rails}

The reviewed design uses an LT8609S step-down regulator for the 5 V
logic rail. The device accepts an input range that includes a fully
charged 6S pack \cite{analog2026lt8609s}. The reviewed schematic used a 2.2 \ensuremath{\mu}H inductor,
1 M\ensuremath{\Omega} and 182 k\ensuremath{\Omega} feedback resistors, a 10 pF feed-forward capacitor, a 1
\ensuremath{\mu}F input capacitor, and a 22 \ensuremath{\mu}F output capacitor, producing a nominal
output near 5.03 V. Those values must be checked against the exact
switching frequency, inductor current rating, capacitor bias derating,
minimum load, loop stability guidance, and the
manufacturer\textquotesingle s recommended layout.

The CAN Pal, ISM330DHCX interface, and display logic were previously
found connected or labeled inconsistently with their intended 3.3 V
operation. The final schematic should explicitly identify the supply
accepted by each installed module rather than assuming that a breakout
labeled VIN is equivalent to the IC supply. The RadioMaster receiver is
assigned to the 5 V rail in the current interface plan. Decoupling must
be placed at each device or connector, and the 3.3 V regulator capacity
must include display and transceiver transient current if those loads
share the controller-board regulator.

\section{Controller, Sensor, Receiver, and Display Interfaces}\label{controller-sensor-receiver-and-display-interfaces}

The active ESP32 pin contract inferred from the current firmware and
schematic review is summarized in Table 5.3. The IMU uses SPI; the
receiver uses a CRSF UART; the display uses a separate SPI-style
interface; and the TWAI controller connects through TXD and RXD to the
external transceiver. The table is a design baseline, not permission to
fabricate. Every pin must be checked against the exact ESP32 board
schematic, boot-strapping constraints, input-only pins, voltage domains,
and the current repository configuration.

\clearpage
\begin{landscape}
\par\medskip\noindent\singlespacing
\phantomsection\addcontentsline{lot}{table}{\protect\numberline{5.3}ESP32 prototype pin assignments}
\textbf{Table 5.3. ESP32 prototype pin assignments}\par

\begin{longtable}[]{@{}
  >{\raggedright\arraybackslash}p{(\columnwidth - 8\tabcolsep) * \real{0.1955}}
  >{\raggedright\arraybackslash}p{(\columnwidth - 8\tabcolsep) * \real{0.2030}}
  >{\raggedright\arraybackslash}p{(\columnwidth - 8\tabcolsep) * \real{0.1203}}
  >{\raggedright\arraybackslash}p{(\columnwidth - 8\tabcolsep) * \real{0.1654}}
  >{\raggedright\arraybackslash}p{(\columnwidth - 8\tabcolsep) * \real{0.3158}}@{}}
\toprule\noalign{}
\begin{minipage}[b]{\linewidth}\raggedright
\textbf{Subsystem}
\end{minipage} & \begin{minipage}[b]{\linewidth}\raggedright
\textbf{Signal}
\end{minipage} & \begin{minipage}[b]{\linewidth}\raggedright
\textbf{ESP32 GPIO}
\end{minipage} & \begin{minipage}[b]{\linewidth}\raggedright
\textbf{Direction at MCU}
\end{minipage} & \begin{minipage}[b]{\linewidth}\raggedright
\textbf{Electrical note}
\end{minipage} \\
\midrule\noalign{}
\endhead
\bottomrule\noalign{}
\endlastfoot
ISM330DHCX & SCLK & 18 & Output & 3.3 V SPI clock \\
ISM330DHCX & SDO\slash{}MISO & 19 & Input & Sensor output to MCU \\
ISM330DHCX & SDI\slash{}MOSI & 23 & Output & MCU output to sensor \\
ISM330DHCX & CS & 27 & Output & Active-low chip select \\
ISM330DHCX & INT1 / INT2 & 34 / 35 & Input & Input-only GPIOs; verify
pull configuration \\
ST7789 display & DIN / CLK & 13 / 14 & Output & 3.3 V logic \\
ST7789 display & CS / DC / RST & 25 / 21 / 22 & Output & Use DC label,
not ambiguous DS \\
RadioMaster RP2 V2 & Receiver TX \ensuremath{\rightarrow} MCU RX & 16 & Input & Verify crossed
UART naming at connector \\
RadioMaster RP2 V2 & MCU TX \ensuremath{\rightarrow} receiver RX & 17 & Output & CRSF
control\slash{}telemetry path \\
CAN transceiver & TXD & 26 & Output & MCU TWAI transmit to
transceiver \\
CAN transceiver & RXD & 32 & Input & Transceiver receive to MCU \\
\end{longtable}
\end{landscape}
\clearpage

\emph{Note.} A September schematic review contained inconsistent TX\slash{}RX
naming in different descriptions. The final connector drawing must name
signals by source and destination and be checked against firmware.

\begin{longtable}[]{@{}
  >{\raggedright\arraybackslash}p{(\columnwidth - 0\tabcolsep) * \real{1.0000}}@{}}
\toprule\noalign{}
\begin{minipage}[b]{\linewidth}\raggedright
\textless INSERT FIGURE 5.3: Insert the finalized controller-interface
schematic with connector pin numbers, voltage rails, decoupling, pull
resistors, ESD protection, test points, and net labels matching Table
5.3. Include board\slash{}module part numbers and footprint
references\textgreater{}\\
What the figure should show: Insert the finalized controller-interface
schematic with connector pin numbers, voltage rails, decoupling, pull
resistors, ESD protection, test points, and net labels matching Table
5.3. Include board\slash{}module part numbers and footprint references\\
Likely source: final KiCad schematic\strut
\end{minipage} \\
\midrule\noalign{}
\endhead
\bottomrule\noalign{}
\endlastfoot
\end{longtable}

\begin{center}\singlespacing
\phantomsection\addcontentsline{lof}{figure}{\protect\numberline{5.3}Controller, sensor, receiver, display, and CAN interface schematic.}
\textbf{Figure 5.3. Controller, sensor, receiver, display, and CAN interface schematic.}
\end{center}

\section{CAN Bus}\label{can-bus}

The current firmware uses the ESP32 TWAI peripheral at 1 Mbit\slash{}s. The
TJA1051T\slash{}3 translates the logic signals to the differential bus
\cite{espressif2026twai}, \cite{nxp2026tja1051}. The main cable should form one linear trunk from one
physical end to the other. Each motor node connects with a short stub,
and CANH\slash{}CANL should be routed as a coupled pair with a nearby reference
conductor. Exactly two 120 \ensuremath{\Omega} terminations belong at the two physical
ends of the assembled bus \cite{ti2026tida01238}. Measuring approximately 60 \ensuremath{\Omega} between
CANH and CANL with all power removed is a useful static check when both
terminators are installed.

The current schematic review identified J9 through J12 as CAN connectors
with pin 1 labeled CANH and pin 2 labeled CANL, but the motor harness
mapping remained ambiguous and must not be inferred from wire color. The
AK40-10 connector and extension harness should be traced end-to-end and
compared with the exact actuator manual before power is applied. CAN
ground or a reference conductor should accompany the differential pair,
while actuator power uses separate conductors sized for motor current.

\begin{longtable}[]{@{}
  >{\raggedright\arraybackslash}p{(\columnwidth - 0\tabcolsep) * \real{1.0000}}@{}}
\toprule\noalign{}
\begin{minipage}[b]{\linewidth}\raggedright
\textless INSERT FIGURE 5.4: Draw the physical CAN trunk from the
controller\slash{}transceiver through all four motor nodes in installed order.
Show cable lengths, stub lengths, CANH, CANL, reference\slash{}ground, shield
if used, node IDs, termination resistor locations, and connector pin
numbers\textgreater{}\\
What the figure should show: Draw the physical CAN trunk from the
controller\slash{}transceiver through all four motor nodes in installed order.
Show cable lengths, stub lengths, CANH, CANL, reference\slash{}ground, shield
if used, node IDs, termination resistor locations, and connector pin
numbers\\
Likely source: final wiring harness drawing and physical
measurement\strut
\end{minipage} \\
\midrule\noalign{}
\endhead
\bottomrule\noalign{}
\endlastfoot
\end{longtable}

\begin{center}\singlespacing
\phantomsection\addcontentsline{lof}{figure}{\protect\numberline{5.4}Recommended linear CAN topology and termination locations.}
\textbf{Figure 5.4. Recommended linear CAN topology and termination locations.}
\end{center}

\section{Printed-Circuit-Board Design}\label{printed-circuit-board-design}

The reviewed repository contained KiCad project, schematic, and PCB
files under 03\_electrical. Earlier PCB evidence showed a
controller-header template with no routed segments, vias, or zones;
later schematic work added the power converter and peripheral
connectors, but manufacturing and validation were not demonstrated. The
final thesis should include the schematic revision, PCB revision, layer
stack, copper weight, design-rule check report, electrical-rule check
report, fabrication outputs, assembly record, and bring-up results.

Signal routing should preserve an uninterrupted reference plane and keep
switch-node copper compact. The LT8609S input capacitor, catch-current
loop, inductor, output capacitors, and feedback divider should follow
the data-sheet placement guidance \cite{analog2026lt8609s}. CAN transient protection
should sit close to the external connector, and the transceiver
decoupling path should be short \cite{nxp2026tja1051}, \cite{ti2026tida01238}. IMU placement should
minimize vibration and thermal gradients while maintaining a known rigid
orientation. High-current actuator paths should be routed as broad
copper areas on both outer layers with dense via stitching and without
shared narrow necks.

\clearpage
\begin{landscape}
\par\medskip\noindent\singlespacing
\phantomsection\addcontentsline{lot}{table}{\protect\numberline{5.4}Recommended printed-circuit-board routing classes}
\textbf{Table 5.4. Recommended printed-circuit-board routing classes}\par

\begin{longtable}[]{@{}
  >{\raggedright\arraybackslash}p{(\columnwidth - 8\tabcolsep) * \real{0.1560}}
  >{\raggedright\arraybackslash}p{(\columnwidth - 8\tabcolsep) * \real{0.1773}}
  >{\raggedright\arraybackslash}p{(\columnwidth - 8\tabcolsep) * \real{0.2128}}
  >{\raggedright\arraybackslash}p{(\columnwidth - 8\tabcolsep) * \real{0.1702}}
  >{\raggedright\arraybackslash}p{(\columnwidth - 8\tabcolsep) * \real{0.2837}}@{}}
\toprule\noalign{}
\begin{minipage}[b]{\linewidth}\raggedright
\textbf{Net class}
\end{minipage} & \begin{minipage}[b]{\linewidth}\raggedright
\textbf{Track or copper width}
\end{minipage} & \begin{minipage}[b]{\linewidth}\raggedright
\textbf{Via pad / drill}
\end{minipage} & \begin{minipage}[b]{\linewidth}\raggedright
\textbf{Clearance}
\end{minipage} & \begin{minipage}[b]{\linewidth}\raggedright
\textbf{Application and notes}
\end{minipage} \\
\midrule\noalign{}
\endhead
\bottomrule\noalign{}
\endlastfoot
Signals & 0.25 mm & 0.60 / 0.30 mm & Fabricator default \ensuremath{\geq} approved
minimum & GPIO, UART, SPI, interrupts, CAN logic; route CANH\slash{}CANL as a
controlled pair where practical. \\
3.3 V & 0.50 mm & 0.80 / 0.40 mm & Approved logic clearance & Peripheral
distribution; use local planes\slash{}pours where helpful. \\
5 V and LT8609 VIN\slash{}VOUT & 1.00 mm & 1.00 / 0.50 mm & Approved power
clearance & Short current loops; verify copper temperature and converter
layout. \\
Each motor branch VBAT and return & 5.00 mm top + bottom pour & 1.00 /
0.50 mm; 8 vias per net at transitions & 0.50 mm minimum; 1.00 mm where
possible & One actuator branch; stitch layers approximately every 5
mm. \\
Each two-motor input branch & 10.00 mm top + bottom pour & 1.00 / 0.50
mm; 16 vias per net at transitions & 0.50 mm minimum; 1.00 mm where
possible & No shared narrow neck; independent input branch for each
motor pair. \\
LT8609 exposed-pad thermal vias & --- & 0.50 / 0.25 mm; five vias & Per
footprint & Confirm against selected package and assembly process. \\
\end{longtable}
\end{landscape}
\clearpage

\emph{Note.} These are design values selected for the present layout,
not universal ampacity guarantees. Verify using the final stackup,
IPC-2152-based analysis, conductor length, allowed temperature rise,
enclosure airflow, and load testing \cite{ipc2009ipc2152}. One-ounce copper is
accepted only with two independent XT60 branches, 10 mm top-and-bottom
pair pours, 5 mm top-and-bottom motor pours, dense stitching, 14 AWG
harness conductors, and no common neck. Two-ounce copper remains
preferable.

\begin{longtable}[]{@{}
  >{\raggedright\arraybackslash}p{(\columnwidth - 0\tabcolsep) * \real{1.0000}}@{}}
\toprule\noalign{}
\begin{minipage}[b]{\linewidth}\raggedright
\textless INSERT FIGURE 5.5: Insert publication-quality exports of the
final schematic sheets. At minimum show power entry and protection,
LT8609S regulator, controller headers, IMU, CAN transceiver and
termination, four motor connectors, receiver, display, debug interface,
test points, and net labels\textgreater{}\\
What the figure should show: Insert publication-quality exports of the
final schematic sheets. At minimum show power entry and protection,
LT8609S regulator, controller headers, IMU, CAN transceiver and
termination, four motor connectors, receiver, display, debug interface,
test points, and net labels\\
Likely source: KiCad schematic plot\strut
\end{minipage} \\
\midrule\noalign{}
\endhead
\bottomrule\noalign{}
\endlastfoot
\end{longtable}

\begin{center}\singlespacing
\phantomsection\addcontentsline{lof}{figure}{\protect\numberline{5.5}Final KiCad schematic sheet set.}
\textbf{Figure 5.5. Final KiCad schematic sheet set.}
\end{center}

\begin{longtable}[]{@{}
  >{\raggedright\arraybackslash}p{(\columnwidth - 0\tabcolsep) * \real{1.0000}}@{}}
\toprule\noalign{}
\begin{minipage}[b]{\linewidth}\raggedright
\textless INSERT FIGURE 5.6: Insert top and bottom PCB renders with
high-current VBAT\slash{}return paths, logic planes, CAN pair, switch node,
regulator loops, IMU placement, connector orientation, fuses, test
points, mounting holes, and thermal vias annotated\textgreater{}\\
What the figure should show: Insert top and bottom PCB renders with
high-current VBAT\slash{}return paths, logic planes, CAN pair, switch node,
regulator loops, IMU placement, connector orientation, fuses, test
points, mounting holes, and thermal vias annotated\\
Likely source: KiCad PCB editor and 3D viewer\strut
\end{minipage} \\
\midrule\noalign{}
\endhead
\bottomrule\noalign{}
\endlastfoot
\end{longtable}

\begin{center}\singlespacing
\phantomsection\addcontentsline{lof}{figure}{\protect\numberline{5.6}Final printed-circuit-board layout with power and signal routing identified.}
\textbf{Figure 5.6. Final printed-circuit-board layout with power and signal routing identified.}
\end{center}

\section{Electrical Safety and Verification}\label{electrical-safety-and-verification}

Initial power-up should use a current-limited bench supply or a fused
battery interface with the actuators disconnected. Resistance-to-ground,
rail polarity, and regulator output should be checked before installing
the controller and modules. Logic rails should then be loaded and
observed for startup overshoot and ripple. The CAN transceiver should be
verified with the motors unpowered before actuator power is enabled.
Each motor branch should be enabled individually while current,
connector temperature, and frame error counters are monitored.

The final electrical design should protect against reverse polarity,
overcurrent, undervoltage, CAN transients, regulator overheating, and
accidental re-arming. The emergency stop must be evaluated as a system
function: removing only a software command is not equivalent to
isolating actuator energy. The implemented behavior, recovery procedure,
and residual powered functions must be documented.

\begin{longtable}[]{@{}
  >{\raggedright\arraybackslash}p{(\columnwidth - 0\tabcolsep) * \real{1.0000}}@{}}
\toprule\noalign{}
\begin{minipage}[b]{\linewidth}\raggedright
\textless VERIFY FROM HARDWARE: Record battery manufacturer and part
number, cell count, capacity, discharge rating, maximum and minimum
operating voltage, connector and wire gauge, fuse types and ratings,
master-disconnect hardware, emergency-stop architecture, BMS\slash{}protection
implementation, converter efficiency and temperatures, rail ripple, peak
voltage droop, and measured actuator branch current.\textgreater{}
\end{minipage} \\
\midrule\noalign{}
\endhead
\bottomrule\noalign{}
\endlastfoot
\end{longtable}
